%%%%%%%%%%%%%%%%%%%%%%%%%%%%%%%%%%%%%%%%%%%%%%%%%%%%%%%%%%%%%%%%%%%%%%%%%%%%%%%
%%  Chapter 17 — Cross-Domain Unification II:
%%               Biochemistry, Linguistics, Semiotics,
%%               Computer Science, and Information Theory
%%
%%  DEPENDENCIES: Chapters 3–16
%%
%%  PURPOSE: Derive biochemistry, language, symbolic systems, computation,
%%           and information theory as operator-algebraic substructures of CIIR.
%%%%%%%%%%%%%%%%%%%%%%%%%%%%%%%%%%%%%%%%%%%%%%%%%%%%%%%%%%%%%%%%%%%%%%%%%%%%%%%

\chapter{Cross-Domain Unification~II: Matter, Language, and Computation}
\label{ch:cross-domain-II}

\vspace{0.5em}
\noindent
This chapter extends the domain-projection apparatus
(Definition~\ref{def:16.2}) to biochemistry, linguistics, semiotics,
computer science, and information theory.  Each domain inherits its
dynamics from the CIIR Lindblad generator and its geometry from the
informational manifold.

% ============================================================
\section{Domain V: Biochemistry and Molecular Biology}
\label{sec:ch17:biochem}
% ============================================================

\subsection{Formal Definition}

\begin{definition}[Molecular State Space]\label{def:17.1}
A \emph{molecular state} is an element $m \in \CS$ (the configuration
space of the primitive triple, Definition~\ref{def:14.1}) equipped with
a finite-dimensional internal Hilbert space $\HC_m$ encoding electronic
and vibrational degrees of freedom.  The \emph{biochemical domain} is:
\[
  D_{\mathrm{chem}} := (P_{\mathrm{chem}},\;
                        \mathcal{O}_{\mathrm{chem}},\;
                        \Omega_{\mathrm{chem}})
\]
where $P_{\mathrm{chem}}$ projects onto the subspace
\[
  \HC_{\mathrm{chem}} := \bigoplus_{m \in \mathcal{M}_{\mathrm{mol}}} \HC_m,
\]
with $\mathcal{M}_{\mathrm{mol}} \subset \CS$ the set of molecular
configurations.
\end{definition}

\subsection{Ontological Mapping}

\begin{center}
\begin{tabular}{lll}
\toprule
\textbf{Biochemical Concept} & \textbf{CIIR Object} & \textbf{Reference} \\
\midrule
Molecule & $m \in \CS$, $\rho_m \in \Omega_{\mathrm{chem}}$ & Def.~\ref{def:17.1} \\
Chemical reaction & $R_{ij} \in \mathcal{O}_{\mathrm{chem}}$ & Def.~\ref{def:17.2} \\
Reaction rate & $\gamma_{ij} \geq 0$ & Def.~\ref{def:17.2} \\
Free energy & $\CC(m)$ (constraint functional) & Def.~\ref{def:14.1} \\
\bottomrule
\end{tabular}
\end{center}

\subsection{Mathematical Construction}

\begin{definition}[Reaction Operator]\label{def:17.2}
A \emph{reaction operator} $R_{ij} : \HC_{m_i} \to \HC_{m_j}$ is a
partial isometry in $\mathcal{O}_{\mathrm{chem}}$ satisfying:
\[
  R_{ij}^\dagger R_{ij} = P_{m_i}, \qquad
  R_{ij} R_{ij}^\dagger = P_{m_j},
\]
where $P_{m_i}$ is the projection onto $\HC_{m_i}$.  The full
reaction dynamics are generated by the Lindblad equation:
\[
  \CL_{\mathrm{chem}}(\rho) = -\frac{i}{\heff}[\hatH_{\mathrm{chem}}, \rho]
  + \sum_{(i,j)} \gamma_{ij}
    \Bigl( R_{ij} \rho R_{ij}^\dagger
    - \tfrac{1}{2}\{R_{ij}^\dagger R_{ij}, \rho\} \Bigr),
\]
where $\hatH_{\mathrm{chem}} := \sum_m \CC(m) P_m$ is the chemical
Hamiltonian constructed from the constraint functional.
\end{definition}

\subsection{Derivation: Thermodynamic Arrow}

\begin{proposition}[Constraint-Induced Free-Energy Minimisation]\label{prop:17.1}
Under $\CL_{\mathrm{chem}}$, the expected constraint decreases
monotonically:
\[
  \frac{d}{dt} \Tr(\hatH_{\mathrm{chem}} \rho)
  = \sum_{(i,j)} \gamma_{ij}\bigl(\CC(m_j) - \CC(m_i)\bigr)
    \Tr(R_{ij}^\dagger R_{ij} \rho) \leq 0,
\]
provided $\gamma_{ij} > 0$ only when $\CC(m_j) \leq \CC(m_i)$
(reactions proceed towards lower constraint).  This recovers the
second law of thermodynamics for chemical systems as a constraint
minimisation principle.
\end{proposition}

\subsection{Cross-Domain Links}

\begin{itemize}
  \item \textbf{To Biology (Sec.~\ref{sec:ch16:biology}):}
        Molecular reactions are the microscopic constituents of the
        biological fitness landscape.  The biological constraint
        $\CC|_{\Omega_{\mathrm{bio}}}$ is the coarse-graining of
        the chemical Hamiltonian $\hatH_{\mathrm{chem}}$.
  \item \textbf{To Emergent Physics (Ch.~14):} The constraint
        functional $\CC$ is the same primitive that generates
        spacetime geometry.  Chemical free energy is a local evaluation
        of the global constraint.
\end{itemize}

% ============================================================
\section{Domain VI: Linguistics and Language Formation}
\label{sec:ch17:linguistics}
% ============================================================

\subsection{Formal Definition}

\begin{definition}[Linguistic Domain]\label{def:17.3}
Let $\Sigma$ be a finite alphabet and $\Sigma^*$ the free monoid over
$\Sigma$.  The \emph{linguistic domain} is:
\[
  D_{\mathrm{ling}} := (L, \mathcal{O}_{\mathrm{ling}}, \Sigma^*)
\]
where $L : \HC \to \ell^2(\Sigma^*)$ is the \emph{symbolic compression
operator}:
\[
  L := \sum_{w \in \Sigma^*} |w\rangle \langle \phi_w |,
\]
with $\{|\phi_w\rangle\}_{w \in \Sigma^*}$ an orthonormal basis of
$\HC$ indexed by strings, and $\mathcal{O}_{\mathrm{ling}} :=
L \CB(\HC) L^\dagger$.
\end{definition}

\subsection{Ontological Mapping}

\begin{center}
\begin{tabular}{lll}
\toprule
\textbf{Linguistic Concept} & \textbf{CIIR Object} & \textbf{Reference} \\
\midrule
Meaning (semantic state) & $\rho \in \mathfrak{D}(\HC)$ & Def.~\ref{def:16.1} \\
Word/symbol & $w \in \Sigma^*$ (basis label) & Def.~\ref{def:17.3} \\
Syntax rule & Algebraic constraint on $\mathcal{O}_{\mathrm{ling}}$ & Thm.~\ref{thm:17.1} \\
Utterance & $L(\rho) \in \ell^2(\Sigma^*)$ & Def.~\ref{def:17.3} \\
\bottomrule
\end{tabular}
\end{center}

\subsection{Mathematical Construction: Syntax as Algebraic Constraint}

\begin{theorem}[Syntactic Algebra]\label{thm:17.1}
Define the \emph{concatenation operator} $C_{w'} : |w\rangle \mapsto
|ww'\rangle$ for $w' \in \Sigma^*$.  The set of grammatically valid
strings is:
\[
  \mathcal{G} := \bigl\{ w \in \Sigma^* :
  P_{\mathrm{gram}} |w\rangle \neq 0 \bigr\},
\]
where $P_{\mathrm{gram}}$ is a projection operator in
$\mathcal{O}_{\mathrm{ling}}$ constructed from the constraint
operator:
\[
  P_{\mathrm{gram}} := \mathbf{1}_{[0,\lambda_{\mathrm{gram}}]}(L \hatC L^\dagger),
\]
with $\lambda_{\mathrm{gram}} > 0$ a threshold below which
constraint is sufficiently low for communication.  Then:
\begin{enumerate}
  \item[\textup{(i)}] $\mathcal{G}$ is a regular language (closed under
    the operations of $P_{\mathrm{gram}}$).
  \item[\textup{(ii)}] The syntactic rules are the algebraic
    constraints defining the range of $P_{\mathrm{gram}}$.
  \item[\textup{(iii)}] Compositional semantics arises from the
    homomorphism property:
    $L(\rho_1 \otimes \rho_2) = L(\rho_1) \star L(\rho_2)$,
    where $\star$ is the convolution on $\ell^2(\Sigma^*)$.
\end{enumerate}
\end{theorem}

\subsection{Cross-Domain Links}

\begin{itemize}
  \item \textbf{To Cognitive Science (Sec.~\ref{sec:ch16:cogsci}):}
        Language production is the application of $L$ to the cognitive
        state: $L \circ \Pi(\rho) \in \ell^2(\Sigma^*)$.  Linguistic
        competence is a property of the domain projection
        $P_{\mathrm{cog}}$.
  \item \textbf{To Computer Science (Sec.~\ref{sec:ch17:cs}):}
        Formal grammars are a special case of the syntactic algebra
        (Theorem~\ref{thm:17.1}), establishing the Chomsky hierarchy
        as a hierarchy of constraint projections.
\end{itemize}

% ============================================================
\section{Domain VII: Semiotics and Symbolic Systems}
\label{sec:ch17:semiotics}
% ============================================================

\subsection{Formal Definition}

\begin{definition}[Semiotic Domain]\label{def:17.4}
The \emph{semiotic domain} generalises the linguistic domain by
replacing the string algebra $\Sigma^*$ with an arbitrary
\emph{sign algebra} $(\mathcal{S}, \circ)$:
\[
  D_{\mathrm{sem}} := (S, \mathcal{O}_{\mathrm{sem}}, \mathcal{S})
\]
where $S : \HC \to \ell^2(\mathcal{S})$ is a \emph{sign operator}
and $\mathcal{S}$ is a monoid of signs with composition $\circ$.
\end{definition}

\begin{definition}[Triadic Sign Structure]\label{def:17.5}
A \emph{CIIR sign} is a triple $(r, o, i)$ where:
\begin{itemize}
  \item $r \in \CR$ is the \emph{referent} (ontic state),
  \item $o \in \mathcal{O}_{\mathrm{sem}}$ is the \emph{sign operator}
        (representamen),
  \item $i := \Phi(o)(r) \in \HC$ is the \emph{interpretant}
        (cognitive state produced by applying the interface map to
        the sign operator acting on the referent).
\end{itemize}
This recovers Peirce's triadic sign model as a special case of the
CIIR interface map.
\end{definition}

\begin{theorem}[Semiotic Functoriality]\label{thm:17.2}
The assignment $r \mapsto (r, o_r, i_r)$ defines a functor
$\mathbf{Sem} : \mathbf{CR} \to \mathbf{Sign}$ from the category
of reality-space objects to the category of signs, where morphisms
are constraint-preserving maps.  This functor factors through the
CIIR representation functor $F$ (Chapter~6).
\end{theorem}

\subsection{Cross-Domain Links}

\begin{itemize}
  \item \textbf{To Linguistics (Sec.~\ref{sec:ch17:linguistics}):}
        The linguistic domain is the sub-category of the semiotic
        domain where the sign algebra is the free monoid
        $\Sigma^*$.
  \item \textbf{To Cognitive Science (Sec.~\ref{sec:ch16:cogsci}):}
        The interpretant $i$ is a cognitive state in
        $\Omega_{\mathrm{cog}}$, establishing semiotics as the
        study of constraint-mediated meaning.
\end{itemize}

% ============================================================
\section{Domain VIII: Computer Science}
\label{sec:ch17:cs}
% ============================================================

\subsection{Formal Definition}

\begin{definition}[Computational Domain]\label{def:17.6}
The \emph{computational domain} is:
\[
  D_{\mathrm{comp}} := (\mathrm{End}(\HC),\;
                        \mathcal{O}_{\mathrm{comp}},\;
                        \Omega_{\mathrm{comp}})
\]
where $\mathcal{O}_{\mathrm{comp}} := \mathrm{End}(\HC)$ is the
full endomorphism algebra and $\Omega_{\mathrm{comp}} :=
\mathfrak{D}(\HC)$.  A \emph{computation} is a finite composition
of operators in $\mathrm{End}(\HC)$.
\end{definition}

\subsection{Theorem: Computation as Endomorphisms}

\begin{theorem}[Turing Completeness]\label{thm:17.3}
Let $\HC$ be infinite-dimensional and separable.  Then:
\[
  \text{Turing-computable functions} \hookrightarrow \mathrm{End}(\HC).
\]
Specifically, for any Turing machine $T$ with tape alphabet $\Gamma$,
there exists a sequence of operators
$U_1, U_2, \ldots \in \mathrm{End}(\HC)$ such that the composition
$U_n \circ \cdots \circ U_1$ applied to an initial state
$|\psi_0\rangle$ encoding the input yields a final state encoding
the output of $T$ on that input (if $T$ halts).

Conversely, every computable element of $\mathrm{End}(\HC)$ (in the
sense of computable operator theory) corresponds to a Turing-computable
function.
\end{theorem}

\begin{proof}[Proof sketch]
Encode the tape configuration as $|\psi\rangle = \sum_{i} c_i |q_i,
s_1 \ldots s_n, h\rangle$ where $q_i$ is the state, $s_j \in \Gamma$,
and $h$ is the head position.  Each transition rule
$(q, s) \to (q', s', \mathrm{dir})$ corresponds to a unitary
$U_{\mathrm{step}}$ acting on this encoding.  The Church--Turing
thesis then gives the embedding.  For the converse, any computable
operator is specified by a recursive function on its matrix elements,
which is Turing-computable by definition.
\end{proof}

\subsection{Derivation: Computational Complexity as Constraint Depth}

\begin{definition}[Constraint Depth]\label{def:17.7}
The \emph{constraint depth} of a computation $U = U_n \circ \cdots
\circ U_1$ is:
\[
  \mathrm{depth}(U) := \sum_{k=1}^{n} \Tr(\hatC\, U_k^\dagger U_k).
\]
This measures the total constraint cost of the computation.
\end{definition}

\begin{proposition}[Complexity--Constraint Correspondence]\label{prop:17.2}
The time complexity $T(n)$ of a Turing machine on input of size $n$
satisfies:
\[
  T(n) \leq \mathrm{depth}(U_n) \cdot \|\hatC^{-1}\|_{\mathrm{op}},
\]
establishing that computational complexity is bounded by constraint
depth.
\end{proposition}

\subsection{Cross-Domain Links}

\begin{itemize}
  \item \textbf{To Linguistics (Sec.~\ref{sec:ch17:linguistics}):}
        Formal grammars are computable operators:
        $P_{\mathrm{gram}} \in \mathrm{End}(\HC)$.  The Chomsky
        hierarchy maps to a hierarchy of constraint depths.
  \item \textbf{To Physics (Ch.~14):} Physical dynamics
        ($e^{-i\hatH t/\heff}$) are unitary elements of
        $\mathrm{End}(\HC)$, establishing physics as a special case
        of computation on the informational manifold.
\end{itemize}

% ============================================================
\section{Domain IX: Information Theory}
\label{sec:ch17:info}
% ============================================================

\subsection{Formal Definition}

\begin{definition}[Information-Theoretic Domain]\label{def:17.8}
The \emph{information-theoretic domain} is:
\[
  D_{\mathrm{info}} := (\mathfrak{D}(\HC),\;
                        S,\;
                        \CD)
\]
where $S(\rho) := -\Tr(\rho \ln \rho)$ is the von Neumann entropy
and $\CD$ is the divergence function (Definition~\ref{def:14.1}).
\end{definition}

\subsection{Theorem: Shannon Theory as Classical Limit}

\begin{theorem}[Classical Information Recovery]\label{thm:17.4}
Let $\rho = \sum_i p_i |i\rangle\langle i|$ be a diagonal density
matrix.  Then:
\begin{enumerate}
  \item[\textup{(i)}] $S(\rho) = H(\{p_i\}) := -\sum_i p_i \ln p_i$
    (Shannon entropy).
  \item[\textup{(ii)}] $\CD(\rho \| \sigma) = D_{\mathrm{KL}}(p \| q)$
    for diagonal $\sigma = \sum_i q_i |i\rangle\langle i|$
    (Kullback--Leibler divergence).
  \item[\textup{(iii)}] The Fisher information metric (Theorem~14.1)
    reduces to the classical Fisher information
    $g_{ij}^F = \sum_k \frac{1}{p_k} \frac{\partial p_k}{\partial
    \theta_i} \frac{\partial p_k}{\partial \theta_j}$.
\end{enumerate}
Thus classical information theory is the diagonal (commutative)
restriction of the CIIR information-geometric framework.
\end{theorem}

\subsection{Derivation: Channel Capacity from Constraint Bandwidth}

\begin{proposition}[Constraint-Bounded Channel Capacity]\label{prop:17.3}
Let the interface map $\Phi$ be viewed as a quantum channel.  Its
Holevo capacity satisfies:
\[
  \chi(\Phi) \leq \ln(\dim \HC_\CM)
  \leq \ln\bigl(\lfloor \kappa_{\max} \cdot \mathrm{vol}(\CM) \rfloor
  + 1\bigr),
\]
relating the channel capacity of the interface to the constraint
curvature and manifold volume.  This recovers the
Shannon--Hartley theorem structure with constraint-geometric
bandwidth.
\end{proposition}

\subsection{Cross-Domain Links}

\begin{itemize}
  \item \textbf{To Computer Science (Sec.~\ref{sec:ch17:cs}):}
        Kolmogorov complexity is the minimum constraint depth
        (Definition~\ref{def:17.7}) of a program producing a given
        state.
  \item \textbf{To Emergent Physics (Ch.~14):} The von Neumann
        entropy $S(\rho)$ is the thermodynamic entropy of the emergent
        physical system.  The information metric is the spacetime
        metric.
\end{itemize}
